\chapter{Geltungsbereich und Grundbegriffe}

Dieser Leitfaden ergänzt Band~9 um eine ausgewählte Übersicht der
wichtigsten mengenbasierten Fassungen und unmittelbaren Folgen des
\emph{Auswahlprinzips}. Vorausgesetzt wird die Mengenlehre in ZF;
das Auswahlaxiom, kurz \textbf{AC}, wird ausdrücklich hinzugenommen.
Die Begriffe Funktion, Familie und Produkt bezeichnen ausschließlich
Mengen und Mengen von geordneten Paaren. Beweise und Beweisskizzen
entfallen. Das Format folgt den Leitfäden zu Band~3, Band~4, Band~5,
Band~28 und Band~37.

\section{Gegenstand und Abgrenzung}

Im Mittelpunkt stehen Auswahlmengen, Auswahlfunktionen, die Existenz
von Produktelementen, die Fortsetzung vorgegebener Auswahlen,
die Auswahl unter Mengenbedingungen sowie vollständige
Distributivgesetze. Die Mengen müssen weder endlich noch abzählbar sein.
Der Leitfaden charakterisiert das Auswahlprinzip und beschreibt,
welche gemeinsamen Auswahlen es ermöglicht.

Allgemeine Regeln der Mengenlehre und Funktionstheorie werden
nicht als Resultate wiederholt. Injektivität, Surjektivität,
Bijektivität und Rechtsinverse bleiben den entsprechenden
Funktionsbänden vorbehalten. Ebenso entfallen die Theorie totaler
Relationen, Ordnungs- und Äquivalenzrelationen, Wohlordnungssatz,
Zornsches Lemma und Anwendungen auf algebraische oder topologische
Strukturen. Die für die Aussagen benötigten Mengendefinitionen
werden dennoch ausdrücklich angegeben.

\textbf{Bedeutung der Auswahlannahme.}
Eine Kennzeichnung \emph{äquivalent zu AC über ZF} bezieht sich
stets auf die Aussage für \emph{alle} Mengen der angegebenen Art.
Sie behauptet nicht, daß jeder einzelne Spezialfall das volle
Auswahlaxiom benötigt. Insbesondere wird die Auswahl aus einer
endlichen Familie nicht als eigenständiges AC-Resultat geführt.
Die bloße Definition einer Auswahlfunktion oder eines Produkts
setzt ebenfalls kein AC voraus.

Die Zusammenstellung ist ein Referenzband wichtiger Aussagen,
keine vollständige Liste aller bekannten oder ableitbaren Sätze.
Die Aussage-Nummern mit Präfix $9H$ gehören ausschließlich zu
diesem Leitfaden.

\section{Mengenfamilien und Auswahlen}

\begin{definition}[Mengenfamilie und indizierte Mengenfamilie]
Eine \emph{Mengenfamilie} $\mathcal A$ ist eine Menge, deren
Elemente die betrachteten Mengen sind.
Eine \emph{indizierte Mengenfamilie} $(A_i)_{i\in I}$ ordnet
jedem Element einer Indexmenge $I$ eine Menge $A_i$ zu.
Verschiedene Indizes dürfen dieselbe Menge bezeichnen.
Die Schreibweise $\bigcup_{i\in I}A_i$ bezeichnet die Vereinigung
aller Familienglieder; für $I=\varnothing$ ist sie leer.
\end{definition}

\begin{definition}[Auswahlfunktion einer Mengenfamilie]
\label{b09guide:def:choice}
Eine \emph{Auswahlfunktion} für $\mathcal A$ ist eine Funktion
\[
 c\colon\mathcal A\longrightarrow\bigcup\mathcal A,
 \qquad c(A)\in A\quad(A\in\mathcal A).
\]
Für eine indizierte Familie $(A_i)_{i\in I}$ heißt eine Funktion
\[
 f\colon I\longrightarrow\bigcup_{i\in I}A_i,
 \qquad f(i)\in A_i\quad(i\in I)
\]
eine \emph{indizierte Auswahlfunktion}.
Bei wiederholten Familiengliedern muß $f$ nicht an jedem zugehörigen
Index denselben Wert wählen; $c$ hat dagegen die Menge selbst
als Argument.
\end{definition}

\begin{definition}[Paarweise disjunkte Familie und Auswahlmenge]
\label{b09guide:def:transversal}
Eine Mengenfamilie $\mathcal A$ heißt \emph{paarweise disjunkt},
wenn für verschiedene $A,B\in\mathcal A$ stets $A\cap B=\varnothing$
gilt. Eine \emph{Auswahlmenge} für $\mathcal A$ ist eine Menge
$T\subseteq\bigcup\mathcal A$ mit
\[
 \forall A\in\mathcal A\;\exists!x\;(x\in T\cap A).
\]
Sie trifft jedes Familienglied in genau einem Element.
Die Existenz solcher Mengen wird im folgenden für
paarweise disjunkte Familien nichtleerer Mengen behauptet.
Für überlappende Familien ist sie kein allgemeines Auswahlprinzip.
\end{definition}

\begin{axiom}[Auswahlprinzip in reiner Mengenform]
\label{b09guide:ax:choice}
Jede paarweise disjunkte Mengenfamilie $\mathcal A$ mit
$\varnothing\notin\mathcal A$ besitzt eine Auswahlmenge:
\[
 \exists T\subseteq\bigcup\mathcal A\quad
 \forall A\in\mathcal A\;\exists!x\;(x\in T\cap A).
\]
Diese Aussage wird mit $\mathrm{AC}$ bezeichnet.
\end{axiom}

Dies ist die Familienform aus Band~3 und der Ausgangspunkt von Band~9.
Die Beschränkung $T\subseteq\bigcup\mathcal A$ entfernt lediglich
Elemente außerhalb der betrachteten Familie.
Für $\mathcal A=\varnothing$ ist $T=\varnothing$ zulässig.

\section{Produkte und Auswahl auf einer festen Menge}

\begin{definition}[Kartesisches Produkt einer Mengenfamilie]
\label{b09guide:def:product}
Für eine indizierte Mengenfamilie $(A_i)_{i\in I}$ sei
\[
 \prod_{i\in I}A_i
 =\left\{f\;\middle|\;
 \begin{gathered}
 f\colon I\longrightarrow\bigcup_{i\in I}A_i,\\
 \forall i\in I\;(f(i)\in A_i)
 \end{gathered}\right\}.
\]
Das Produkt ist somit die Menge aller indizierten Auswahlfunktionen.
Die Definition schließt
$\prod_{i\in\varnothing}A_i=\{\varnothing\}$ ein:
Der leere Funktionsgraph ist das einzige leere Tupel.
Ein einzelner leerer Faktor bei nichtleerem $I$ schließt
Produktelemente aus.
\end{definition}

\begin{definition}[Auswahl auf allen nichtleeren Teilmengen]
\label{b09guide:def:selector}
Für eine Menge $U$ sei $\mathcal P(U)$ ihre Potenzmenge und
\[
 \ChoiceSet(U)=
 \left\{c\;\middle|\;
 \begin{gathered}
 c\colon\mathcal P(U)\setminus\{\varnothing\}\longrightarrow U,\\
 \forall A\in\mathcal P(U)\setminus\{\varnothing\}\;(c(A)\in A)
 \end{gathered}\right\}.
\]
Ein Element von $\ChoiceSet(U)$ heißt \emph{Auswahlfunktion auf
den nichtleeren Teilmengen von $U$}.
Dabei ist $\ChoiceSet(\varnothing)=\{\varnothing\}$.
\end{definition}

\section{Wegweiser}

Kapitel~\ref{b09guide:ch:criteria} sammelt die äquivalenten
Auswahl- und Produktfassungen.
Kapitel~\ref{b09guide:ch:selections} behandelt Vorgaben,
simultane Auswahlen und die disjunkte Verfeinerung von Überdeckungen.
Kapitel~\ref{b09guide:ch:distribution} enthält die vollständigen
mengentheoretischen Distributivgesetze.
Alle folgenden Sätze setzen AC voraus; ihre jeweils
allgemeine Fassung ist bereits äquivalent zu AC über ZF.
